\chapter{HASIL DAN PEMBAHASAN}

\petunjuk{Pisahkan penyajian hasil dari interpretasinya. Setiap tabel, gambar, dan nilai numerik harus dijelaskan dalam narasi serta dihubungkan dengan tujuan dan kajian pustaka.}

\section{Hasil Implementasi}

Jelaskan hasil implementasi dengan urutan yang konsisten terhadap tahapan metodologi. Gunakan tangkapan layar hanya jika elemen antarmuka perlu dibahas. Pastikan teks di dalam gambar tetap terbaca.

\section{Hasil Pengujian}

Tabel~\ref{tab:hasil-pengujian} menunjukkan contoh penyajian hasil beberapa skenario. Angka pada tabel ini hanya merupakan data contoh dan wajib diganti dengan hasil pengujian yang sebenarnya.

\begin{table}[H]
  \centering
  \caption{Contoh hasil pengujian sistem}
  \label{tab:hasil-pengujian}
  \begin{tabularx}{\textwidth}{YCC}
    \toprule
    \textbf{Skenario} & \textbf{Hasil yang Diharapkan} & \textbf{Status} \\
    \midrule
    Pengguna memasukkan data valid & Data tersimpan & Berhasil \\
    Pengguna memasukkan data kosong & Validasi ditampilkan & Berhasil \\
    Pengguna meminta laporan & Laporan dihasilkan & Berhasil \\
    \bottomrule
  \end{tabularx}
\end{table}

\section{Pembahasan}

Pembahasan tidak hanya mengulang angka hasil pengujian. Jelaskan penyebab hasil, keterkaitannya dengan teori atau penelitian sebelumnya, implikasi praktis, serta kondisi yang membatasi generalisasi hasil.

Ketika membandingkan hasil dengan penelitian lain, gunakan sitasi dari database referensi. Sebagai contoh, konsep penyusunan dokumen teknis terstruktur dapat dirujuk melalui \cite{lamport1994}.
